\section{Trust and reassurance: why this technology can be trusted}
\label{sec:trust}

Medical AI bills face honest public skepticism, so the seventh appeal is
reassurance. Trust is not asserted; it is built from named supports: physician
endorsement, academic review, mandatory safety testing, transparency provisions,
pilot results, and an audit trail that anyone can inspect. The objective is a
committee that can tell constituents, with evidence, that this technology can be
trusted. \autoref{fig:trust-scaffold} names the supports that hold that trust up.

\begin{narrativefig}
\node[nfact] (r) at (0,0) {Trust in medical AI};
\node[nftwo] (a) at (4.2,1.5) {Physician endorsement};
\node[nftwo] (b) at (4.2,0.0) {Independent review};
\node[nftwo] (c) at (4.2,-1.5) {Ten VVUQ gates};
\node[nfhope] (d) at (8.4,1.5) {Transparency};
\node[nfhope] (e) at (8.4,0.0) {Pilot results};
\node[nfhope] (f) at (8.4,-1.5) {Audit trail};
\draw[nfedge] (r)--(a); \draw[nfedge] (r)--(b); \draw[nfedge] (r)--(c);
\draw[nfedge] (a)--(d); \draw[nfedge] (b)--(e); \draw[nfedge] (c)--(f);
\end{narrativefig}
\vspace{-0.15cm}
\figcaption{Figure 9. Trust scaffolding (Mermaid 10 as colored TikZ): confidence
rests on named, inspectable supports, not on assurance alone.}

The strongest support is the most concrete. Each of the ten gates is bound to a
published consensus standard already used in real engineering practice, such as
ASME V\&V 40-2018 for assessing the credibility of computational models in medical
devices \cite{asme2018vv40}, and the surgical-humanoid run cleared all of them,
passing 172 of 172 automated tests \cite{Kawchak2026MobilePancreaticCancerUnitree,
Kawchak2026VVUQOncologyClinicalTrial}. \autoref{fig:trust-gates} shows the funnel,
and \autoref{tab:trust} binds representative gates to their standards and
thresholds, which is what makes the assurance argument traceable to a regulator.

\begin{narrativefig}
\node[nfone] (i) at (0,0) {Candidate action};
\node[nftwo] (g1) at (3.4,0) {Verification fraction 1.0};
\node[nftwo] (g2) at (6.8,0) {Validation agreement};
\node[nfthree] (g3) at (10.2,0) {Standards binding};
\node[nfdec] (g4) at (5.1,-2.0) {Catastrophe predicate};
\node[nfhope] (o) at (9.0,-2.0) {ACCEPT under audit};
\draw[nfedge] (i)--(g1); \draw[nfedge] (g1)--(g2); \draw[nfedge] (g2)--(g3);
\draw[nfedge] (g3) to[out=-90,in=0] (g4.east);
\draw[nfedge] (g4)--(o);
\end{narrativefig}
\vspace{-0.15cm}
\figcaption{Figure 10. The ten-gate funnel (Mermaid 07 as colored TikZ): each gate
is a published threshold, and only a clean pass reaches ACCEPT.}

\begin{table}[H]
\footnotesize
\begin{tabularx}{\textwidth}{@{}L{3.4cm} Y L{3.0cm}@{}}
\toprule
\textbf{Gate} & \textbf{Bound standard or basis} & \textbf{Pass threshold}\\
\midrule
Verification fraction & ASME V\&V 40-2018 credibility & exactly 1.0\\
Validation agreement & Reference cohort comparison & up to 1.00\\
Relative error & Numerical tolerance & as tight as 0.01\\
Catastrophe predicate & Hard safety rule & block before execution\\
\bottomrule
\end{tabularx}
\caption{Table 7. Representative gates bound to standards and thresholds.}
\label{tab:trust}
\end{table}

Reassurance works when it is specific. A named standard and a published test
result persuade where a promise of safety cannot. The quiet question follows. If
every check can be bound to a standard and shown to pass, why would human review
alone remain the only gate on increasingly complex AI outputs?
